\question[5]
Ein Getränk kostet 1,50\,\texteuro. Der Getränkeautomat nimmt nur 50-Cent-Münzen (Eingabe $50$) und
1-Euro-Münzen (Eingabe $100$) an. Sobald genug Geld eingeworfen wurde, gibt er das Getränk und
eventuell Wechselgeld aus.
\begin{parts}
\part Zeichne das Zustandsdiagramm eines Mealy-Automaten.
\part Gib die Zustandsübergangstabelle an.
\part Welche Ausgaben erzeugt der Automat bei der Eingabefolge $50, 100, 100, 100$?
\end{parts}

\begin{solutionbox}{9cm}
\begin{parts}
\part Der Zustand ist der bisher eingeworfene Betrag. -- bedeutet: keine Ausgabe.

\includegraphics[height=3.5cm]{\pfad/mealy_01/mealy_01.png}
\part
\begin{tabular}[t]{r|cc}
 & 50 & 100 \\ \hline
$\rightarrow$ 0\,ct & 50\,ct / -- & 100\,ct / -- \\
50\,ct & 100\,ct / -- & 0\,ct / Getränk \\
100\,ct & 0\,ct / Getränk & 0\,ct / Getränk, 50\,ct zurück \\
\end{tabular}
\part --, Getränk, --, Getränk und 50\,ct zurück.
\end{parts}
\end{solutionbox}
